% ── 5.1  Experimental Setup and Descriptive Statistics ───────────────────
% PURPOSE: Orient the reader before any hypothesis test -- recap the run matrix and
%          report per-policy descriptives so later inferential moves have a baseline.
% INTEGRITY: all numeric descriptives are from the trustworthy 144 rerun
% (runs_144_seq_cceb325, deepseek-v4-flash, sha cceb325; deep audit = 0 flags). The old
% runs_k27_merged values are NOT used here -- the rerun exists because that matrix
% was instrument-crippled (see Section~\ref{sec:instrument}).
% ─────────────────────────────────────────────────────────────────────────

\section{Experimental Setup and Descriptive Statistics}

The evaluation rests on a single balanced matrix of \textbf{144 controlled runs}:
6 memory policies $\times$ 8 official SWE-Bench-CL sequences $\times$ 3 seeds, with
24 runs per policy and 48 runs per seed. Retrieval scoring is held constant across all
six conditions (pure cosine, identical $k$ and context budget), so the only factor that
varies between conditions is \emph{what each policy chooses to keep or forget}. A single
frozen model, \texttt{deepseek-v4-flash}, serves every model role --- the coding agent,
the memory-type classifier, reflection, and CLS consolidation --- so that the model is a
fixed factor across all conditions; between-policy contrasts are therefore valid even
though absolute resolution rates are not comparable to public SWE-Bench leaderboards. The
full slate of model, embedder, cost-metric, and host substitutions is disclosed in
Section~\ref{sec:deviations}, and the instrument-health validation that precedes the
matrix is reported in Section~\ref{sec:instrument}.

All inferential statistics in the sections that follow take the \emph{sequence} as the
unit of analysis ($N=8$ sequence-level means per policy), never the individual task ---
a deliberate guard against pseudo-replication. At the level of individual task executions
the matrix comprises \textbf{4{,}914 task rows}, of which \textbf{2{,}175 were resolved},
for an overall resolution rate of \textbf{44.3\%}; a full deep audit of the trustworthy
rerun (\texttt{runs\_144\_seq\_cceb325}) returned zero integrity flags.

\paragraph{Descriptive picture.}
Table~\ref{tab:descriptives} reports per-policy means over the eight sequences on the
three outcome axes the thesis cares about: resolution rate (the CL-F1 proxy; see the
caveat below), memory footprint carried per task, and total token cost per run. Two
features of the descriptives are visible before any test is run. First, the policies bunch
tightly on resolution --- every condition lands between 0.427 and 0.449, a span of barely
two percentage points --- while differing more than two-fold on footprint, from
\texttt{no\_memory}'s zero through the pruners' \textasciitilde 600 tokens/task up to
\texttt{full\_memory}'s 1{,}247. Second, the tool-call budget is essentially flat across
conditions (20.9--21.1 calls/task), a first hint --- developed in Section~\ref{sec:h4} ---
that the memory policy does not change the agent's step-to-step behaviour. The inferential
sections test whether the small resolution spread is signal or noise; the descriptives
already foreshadow the answer.

\begin{table}[htbp]
  \centering
  \caption{Per-policy descriptive statistics, mean over the eight SWE-Bench-CL
  sequences ($N=8$). Resolution rate is the reported CL-F1 proxy. Footprint is memory
  tokens carried per task; total cost is tokens per run. Rows are ordered by resolution
  rate. Source: \texttt{runs\_144\_seq\_cceb325} (\texttt{deepseek-v4-flash}).}
  \label{tab:descriptives}
  \begin{tabular}{lrrrr}
    \toprule
    Policy & Resolve rate & Footprint & Total tok & Tool calls \\
           & (CL-F1 proxy) & (tok/task) & (M/unit)  & /task      \\
    \midrule
    cls\_consolidation & 0.449 & 704   & 6.08 & 21.11 \\
    random\_prune      & 0.447 & 601   & 6.06 & 21.00 \\
    full\_memory       & 0.445 & 1{,}247 & 6.09 & 21.05 \\
    recency\_prune     & 0.437 & 604   & 6.05 & 20.93 \\
    no\_memory         & 0.429 & 0     & 5.74 & 20.98 \\
    type\_aware\_decay & 0.427 & 598   & 5.96 & 20.98 \\
    \bottomrule
  \end{tabular}
\end{table}

\paragraph{Headroom for a memory effect.}
A precondition for any between-policy contrast to be interpretable is that the experiment
sits clear of both floor and ceiling: if the agent solved almost everything or almost
nothing regardless of memory, no contrast would have outcomes left to move. The overall
resolution rate of \textbf{44.3\%} sits squarely in the mid-range --- far from both the
floor and the ceiling --- so there is ample room for a memory effect to register in either
direction. A near-null result on this matrix is therefore a finding about memory, not an
artefact of a saturated or floored benchmark.

\paragraph{What ``CL-F1'' means here.}
Throughout this chapter the reported CL-F1 equals the per-condition resolution rate
($\text{mean\_cl\_f1} \equiv \text{mean\_resolved\_rate}$); the full anchor-probe
instrumentation that would isolate the CL-specific \emph{stability} dimension was not run
at matrix scale. The stability dimension is instead addressed through a targeted
construct-validity sub-study (Section~\ref{sec:h5}) and disclosed in Threats to Validity
(Section~\ref{sec:threats}). Readers should therefore read CL-F1 in this thesis as a
resolution-rate proxy, not as a forgetting-curve measurement.
